\chapter*{A note on sources and numbers}
\addcontentsline{toc}{chapter}{A note on sources and numbers}
\markboth{A note on sources}{A note on sources}

Cuban statistics are a problem with three layers, and a book that quotes them
without saying so is misleading its reader by arithmetic rather than by
argument.

The first layer is ordinary. Cuba's national statistics office, the Oficina
Nacional de Estad\'istica e Informaci\'on (ONEI), publishes an annual
\emph{Anuario Estad\'istico} that is detailed, internally consistent, and by
the standards of the region reasonably competent. Health and education series
in particular are collected by an administrative apparatus that reaches every
municipality, and the vital-registration coverage is close to complete, which
is more than can be said for most countries at Cuba's income level. Where this
book quotes life expectancy, infant mortality, school enrolment or physician
counts, ONEI is usually the source and the figure is usually trustworthy in the
narrow sense that the state measured what it says it measured.

The second layer is the exchange rate, and it corrupts everything monetary. For
most of the period covered here the Cuban state carried an official rate of one
peso to one dollar for enterprise accounting while households changed money at
twenty-four or twenty-five, and after 2021 at rates that in the informal market
passed three hundred. Any figure denominated in Cuban pesos and converted at an
official rate --- GDP, GDP per capita, wages, the value of the social wage, the
cost of a subsidy --- is therefore not a measurement but an accounting
convention, and cross-country comparisons built on it are close to meaningless.
Cuba's reported GDP per capita has for decades placed it in the upper-middle
band of Latin America while its consumption basket has looked like the lower.
Both facts are real; the contradiction is in the conversion. This book gives
peso figures with the rate attached, prefers physical quantities to monetary
ones wherever a physical quantity exists --- tons of sugar, barrels of oil,
megawatts, doctors per thousand, calories per day --- and treats any single
figure for Cuban GDP as an argument rather than a datum.

The third layer is political, and it runs in both directions. The state does
not publish, or publishes late and partially, series that embarrass it: the
scale of the 1993--94 contraction was reconstructed years afterwards, the
prison population is not disclosed, the count of political detentions is
produced by opposition monitors and by no one else, and after 2020 several
routine series in the \emph{Anuario} simply stopped appearing. Against this,
the exile literature has its own inflations, some of them very large and very
durable --- the twenty thousand dead attributed to Batista's last years, a
figure that originated in a newspaper and has never been supported by anything
resembling a count, is the best-known example, and it is repeated to this day
by people who would never accept a Cuban government figure on the same
evidence. Where a contested number appears in this book it appears as a range
with the disagreement named, and where the honest answer is that nobody knows,
the text says nobody knows.

Three practical conventions follow.

Figures asserted from a secondary source that has not been traced to a primary
one are marked in the manuscript and listed in \texttt{verify.md} in the book's
source directory. The mark prints nothing, so it does not clutter the page; it
exists so that the list of unverified claims is mechanically extractable rather
than a matter of memory.

Where a range is given, the endpoints are the two most defensible published
estimates and not the extremes of what has been claimed. A figure that has been
claimed only by a party with an interest in it, and never independently, is
described that way rather than averaged into a range as though it were
evidence.

Where the book states a causal claim --- that the embargo cost Cuba this much,
that the 2021 monetary reform caused that much inflation --- it states what
would have to be true for the claim to hold, because on this subject almost
every causal claim is contested by someone whose objection is not frivolous.
The two standing questions, how much of Cuba's poverty is the embargo and how
much is the model, and how much of Cuba's social achievement is socialism and
how much is what any determined state can buy with a large enough external
subsidy, are not settled in this book. They are, however, kept visible, because
Part III's design has to work under either answer.
